\vspace{-1.0em}
\paragraph{Hierarchical and Semantic Representations in Language Models.} 
A large body of work has shown that transformer-based language models develop hierarchical internal representations in which different layers encode progressively abstract linguistic information. Early analyses revealed that lower layers primarily capture lexical and syntactic structure, while deeper layers increasingly encode semantic properties of the input. Ethayarajh~\cite{ethayarajh-2019-contextual} observed that contextualized embeddings become progressively anisotropic and semantically structured in higher layers, suggesting increasing abstraction along model depth. Similarly, Tenney et al.~\cite{tenney-etal-2019-bert} demonstrated that BERT implicitly recovers a classical NLP pipeline, with successive layers specializing in part-of-speech tagging, syntactic relations, and semantic reasoning. Related probing studies further showed that linguistic structure such as syntax and dependency trees can be recovered from transformer representations~\cite{hewitt-manning-2019,rogers-etal-2020}.

These findings suggest that semantic stability is not imposed explicitly but emerges naturally through representation learning. As information propagates through the network, paraphrased inputs often converge toward similar internal representations despite large lexical variation. Representational similarity analyses further support this observation, showing that semantically equivalent inputs tend to occupy nearby regions in latent space~\cite{bommasani-etal-2020-interpreting}. While these studies provide empirical evidence that transformer representations encode semantic structure, they do not characterize where such invariance emerges across layers or how it can be systematically quantified.

\vspace{-1.0em}
\paragraph{Measuring Representation Similarity and Invariance. }
Understanding internal representations requires methods for comparing activation structures across inputs, layers, and models. A common approach is representational similarity analysis (RSA), which measures relationships between representation spaces using statistical or geometric metrics. Widely used measures include cosine similarity, canonical correlation analysis (CCA), projection-weighted CCA (PWCCA), and centered kernel alignment (CKA), as well as related approaches such as singular vector CCA (SVCCA) for comparing neural representations across networks~\cite{raghu2017svcca}.

Cosine similarity is frequently used to analyze semantic similarity in embedding spaces and has been applied extensively to contextualized language representations~\cite{ethayarajh-2019-contextual,bommasani-etal-2020-interpreting}. CCA-based approaches compare representations by identifying maximally correlated projections between activation spaces, allowing similarity to be measured even when neuron-wise correspondence is absent. Morcos et al.~\cite{morcos2018pwcca} proposed projection-weighted CCA to emphasize directions that contribute most strongly to representation variance. Centered kernel alignment (CKA) was later introduced as a similarity metric invariant to orthogonal transformations and isotropic scaling, enabling robust comparison of deep neural representations across layers and training runs~\cite{kornblith2019similarity}. These techniques have also been applied to study how representations evolve across transformer depth and across independently trained models~\cite{reif2019visualizing}.

These techniques have been widely used to study representation dynamics in neural networks. However, most prior analyses focus on pairwise similarity between representations rather than explicitly characterizing semantic invariance under controlled transformations. As a result, existing work provides useful tools for representation comparison but offers limited insight into where invariant semantic structures emerge within large language models.

\vspace{-1.0em}
\paragraph{Model-Specific Latent Structure and Transferability.} 
While semantic robustness is often observed within individual models, growing evidence suggests that internal representation structures are strongly model-specific. Parameter-efficient prompt tuning studies show that learned prompt embeddings frequently fail to transfer across models even when architectures are closely related. Lester et al.~\cite{lester-etal-2021-power} demonstrated that soft prompts optimized for one model activate latent pathways that do not generalize to others, indicating that semantic representations are embedded within model-specific internal coordinate systems.

Similar observations arise in adversarial trigger research. Wallace et al.~\cite{wallace-etal-2019-universal} introduced universal adversarial triggers capable of inducing targeted behaviors in language models, yet these triggers rarely generalize across architectures or training runs. Such failures of cross-model transfer suggest that although models may exhibit similar external behavior, their internal activation geometries can differ substantially. This phenomenon reflects the non-convex nature of neural network optimization, where models trained on similar data can converge to distinct representational embeddings of the same underlying semantic structure.

\vspace{-1.0em}
\paragraph{Internal Activations for Model Identification.}
Recent work has begun to leverage internal activations for model identification and forensic analysis. Xu et al.~\cite{xu-etal-2024-instructional} showed that instruction tuning leaves persistent activation signatures in intermediate layers, enabling reliable identification of fine-tuned model variants. These results demonstrate that training procedures can imprint distinctive patterns in latent space that remain observable through activation analysis.

However, existing approaches to model identification typically rely on task-specific signals or training-induced artifacts. In contrast, the present work investigates invariant latent structures that arise naturally from representation learning itself. By analyzing how semantically equivalent inputs organize within latent space and how this organization evolves across depth, we aim to identify zones where semantic invariance emerges as an intrinsic property of the model. This perspective provides a task-agnostic framework that connects representation geometry, semantic stability, and model-specific structure.